\chapter{Chapter06 — Regularity}
%%% generated by the blueprint pipeline from FABL.Chapter06.Pseudorandomness.Regularity
\section{Overview}

Declarations of the Lean module \texttt{FABL.Chapter06.Pseudorandomness.Regularity}.

\section{Declarations}

\begin{definition}[Fourier regularity]
\lean{FABL.IsFourierRegular}
\statementsource{Analysis-of-Boolean-Functions-by-Ryan-ODonnell}{
  pdf-ffbe9c7b3306-p144-b003,
  pdf-ffbe9c7b3306-p169-b008
}
\label{FABL.IsFourierRegular}
\leanok
\uses{FABL.fourierCoeff}
For a real number $\varepsilon$ and a real-valued function $f$ on the sign cube of
length $n$, say that $f$ is Fourier regular at level $\varepsilon$ if every nonempty
subset $S$ of the coordinate indices satisfies $|\widehat f(S)|\le\varepsilon$, where
$\widehat f(S)$ is the Fourier coefficient of $f$ at $S$.
\end{definition}

\begin{definition}[Low-degree Fourier regularity]
\lean{FABL.IsLowDegreeFourierRegular}
\statementsource{Analysis-of-Boolean-Functions-by-Ryan-ODonnell}{
  pdf-ffbe9c7b3306-p069-b003,
  pdf-ffbe9c7b3306-p146-b004,
  pdf-ffbe9c7b3306-p169-b008
}
\label{FABL.IsLowDegreeFourierRegular}
\leanok
\uses{FABL.fourierCoeff}
For a real number $\varepsilon$, a natural number $k$, and a real-valued function $f$ on
the sign cube of length $n$, say that $f$ is Fourier regular through degree $k$ at level
$\varepsilon$ if every nonempty subset $S$ of the coordinate indices with $|S|\le k$
satisfies $|\widehat f(S)|\le\varepsilon$.
\end{definition}

\begin{theorem}[Regularity at bounded degree]
\lean{FABL.IsFourierRegular.isLowDegreeFourierRegular}
\proofsource{Analysis-of-Boolean-Functions-by-Ryan-ODonnell}{pdf-ffbe9c7b3306-p146-b004}
\label{FABL.IsFourierRegular.isLowDegreeFourierRegular}
\leanok
\uses{FABL.IsFourierRegular, FABL.IsLowDegreeFourierRegular}
\difficulty{1}
Let $n$ be a nonnegative integer, let $f$ be a real-valued function on the sign cube
$\{−1,1\}^n$, and let $\varepsilon$ be real. Suppose every nonempty set of
coordinates has Fourier coefficient of absolute value at most $\varepsilon$. Then, for
every nonnegative integer $k$, every nonempty set of at most $k$ coordinates has Fourier
coefficient of absolute value at most $\varepsilon$.
\end{theorem}

\begin{theorem}[Fourier regularity at full dimension]
\lean{FABL.isLowDegreeFourierRegular_dimension_iff}
\proofsource{Analysis-of-Boolean-Functions-by-Ryan-ODonnell}{pdf-ffbe9c7b3306-p146-b004}
\label{FABL.isLowDegreeFourierRegular_dimension_iff}
\leanok
\uses{FABL.IsFourierRegular, FABL.IsFourierRegular.isLowDegreeFourierRegular, FABL.IsLowDegreeFourierRegular}
\difficulty{3}
Let $\varepsilon\in\mathbb R$ and let $f:\{-1,1\}^n\to\mathbb R$. Every nonempty
coordinate set of size at most $n$ has Fourier coefficient of absolute value at most
$\varepsilon$ if and only if every nonempty coordinate set has Fourier coefficient of
absolute value at most $\varepsilon$.
\end{theorem}

\begin{theorem}[Fourier regularity from the uniform $L^1$ norm]
\lean{FABL.isFourierRegular_uniformLpNorm_one}
\label{FABL.isFourierRegular_uniformLpNorm_one}
\leanok
\uses{FABL.IsFourierRegular, FABL.abs_fourierCoeff_le_uniformLpNorm_one, FABL.uniformLpNorm}
\difficulty{1}
Let $n\ge 0$ and let $f:\{\pm1\}^n\to\mathbb R$. Every nonempty-set Fourier coefficient
of $f$ has absolute value at most the uniform average of $|f|$ on the cube.
\end{theorem}

\begin{theorem}[Zero Fourier regularity and constant functions]
\lean{FABL.isFourierRegular_zero_iff_exists_const}
\label{FABL.isFourierRegular_zero_iff_exists_const}
\leanok
\uses{FABL.IsFourierRegular, FABL.expect_monomial, FABL.fourierCoeff, FABL.fourier_expansion, FABL.monomial}
\difficulty{5}
Let $n\ge 0$ and let $f:\{\pm1\}^n\to\mathbb R$. All nonempty-set Fourier coefficients
of $f$ vanish if and only if $f$ is constant on the cube.
\end{theorem}
